\IfFileExists{stacks-project.cls}{%
\documentclass{stacks-project}
}{%
\documentclass{amsart}
}

% For dealing with references we use the comment environment
\usepackage{verbatim}
\newenvironment{reference}{\comment}{\endcomment}
%\newenvironment{reference}{}{}
\newenvironment{slogan}{\comment}{\endcomment}
\newenvironment{history}{\comment}{\endcomment}

% For commutative diagrams we use Xy-pic
\usepackage[all]{xy}

% We use 2cell for 2-commutative diagrams.
\xyoption{2cell}
\UseAllTwocells

% We use multicol for the list of chapters between chapters
\usepackage{multicol}

% This is generally recommended for better output
\usepackage{lmodern}
\usepackage[T1]{fontenc}

% For cross-file-references
\usepackage{xr-hyper}

% Package for hypertext links:
\usepackage{hyperref}

% For any local file, say "hello.tex" you want to link to please
% use \externaldocument[hello-]{hello}
\externaldocument[introduction-]{introduction}
\externaldocument[conventions-]{conventions}
\externaldocument[sets-]{sets}
\externaldocument[categories-]{categories}
\externaldocument[topology-]{topology}
\externaldocument[sheaves-]{sheaves}
\externaldocument[sites-]{sites}
\externaldocument[stacks-]{stacks}
\externaldocument[fields-]{fields}
\externaldocument[algebra-]{algebra}
\externaldocument[brauer-]{brauer}
\externaldocument[homology-]{homology}
\externaldocument[derived-]{derived}
\externaldocument[simplicial-]{simplicial}
\externaldocument[more-algebra-]{more-algebra}
\externaldocument[smoothing-]{smoothing}
\externaldocument[modules-]{modules}
\externaldocument[sites-modules-]{sites-modules}
\externaldocument[injectives-]{injectives}
\externaldocument[cohomology-]{cohomology}
\externaldocument[sites-cohomology-]{sites-cohomology}
\externaldocument[dga-]{dga}
\externaldocument[dpa-]{dpa}
\externaldocument[sdga-]{sdga}
\externaldocument[hypercovering-]{hypercovering}
\externaldocument[schemes-]{schemes}
\externaldocument[constructions-]{constructions}
\externaldocument[properties-]{properties}
\externaldocument[morphisms-]{morphisms}
\externaldocument[coherent-]{coherent}
\externaldocument[divisors-]{divisors}
\externaldocument[limits-]{limits}
\externaldocument[varieties-]{varieties}
\externaldocument[topologies-]{topologies}
\externaldocument[descent-]{descent}
\externaldocument[perfect-]{perfect}
\externaldocument[more-morphisms-]{more-morphisms}
\externaldocument[flat-]{flat}
\externaldocument[groupoids-]{groupoids}
\externaldocument[more-groupoids-]{more-groupoids}
\externaldocument[etale-]{etale}
\externaldocument[chow-]{chow}
\externaldocument[intersection-]{intersection}
\externaldocument[pic-]{pic}
\externaldocument[weil-]{weil}
\externaldocument[adequate-]{adequate}
\externaldocument[dualizing-]{dualizing}
\externaldocument[duality-]{duality}
\externaldocument[discriminant-]{discriminant}
\externaldocument[derham-]{derham}
\externaldocument[local-cohomology-]{local-cohomology}
\externaldocument[algebraization-]{algebraization}
\externaldocument[curves-]{curves}
\externaldocument[resolve-]{resolve}
\externaldocument[models-]{models}
\externaldocument[functors-]{functors}
\externaldocument[equiv-]{equiv}
\externaldocument[pione-]{pione}
\externaldocument[etale-cohomology-]{etale-cohomology}
\externaldocument[proetale-]{proetale}
\externaldocument[relative-cycles-]{relative-cycles}
\externaldocument[more-etale-]{more-etale}
\externaldocument[trace-]{trace}
\externaldocument[crystalline-]{crystalline}
\externaldocument[spaces-]{spaces}
\externaldocument[spaces-properties-]{spaces-properties}
\externaldocument[spaces-morphisms-]{spaces-morphisms}
\externaldocument[decent-spaces-]{decent-spaces}
\externaldocument[spaces-cohomology-]{spaces-cohomology}
\externaldocument[spaces-limits-]{spaces-limits}
\externaldocument[spaces-divisors-]{spaces-divisors}
\externaldocument[spaces-over-fields-]{spaces-over-fields}
\externaldocument[spaces-topologies-]{spaces-topologies}
\externaldocument[spaces-descent-]{spaces-descent}
\externaldocument[spaces-perfect-]{spaces-perfect}
\externaldocument[spaces-more-morphisms-]{spaces-more-morphisms}
\externaldocument[spaces-flat-]{spaces-flat}
\externaldocument[spaces-groupoids-]{spaces-groupoids}
\externaldocument[spaces-more-groupoids-]{spaces-more-groupoids}
\externaldocument[bootstrap-]{bootstrap}
\externaldocument[spaces-pushouts-]{spaces-pushouts}
\externaldocument[spaces-chow-]{spaces-chow}
\externaldocument[groupoids-quotients-]{groupoids-quotients}
\externaldocument[spaces-more-cohomology-]{spaces-more-cohomology}
\externaldocument[spaces-simplicial-]{spaces-simplicial}
\externaldocument[spaces-duality-]{spaces-duality}
\externaldocument[formal-spaces-]{formal-spaces}
\externaldocument[restricted-]{restricted}
\externaldocument[spaces-resolve-]{spaces-resolve}
\externaldocument[formal-defos-]{formal-defos}
\externaldocument[defos-]{defos}
\externaldocument[cotangent-]{cotangent}
\externaldocument[examples-defos-]{examples-defos}
\externaldocument[algebraic-]{algebraic}
\externaldocument[examples-stacks-]{examples-stacks}
\externaldocument[stacks-sheaves-]{stacks-sheaves}
\externaldocument[criteria-]{criteria}
\externaldocument[artin-]{artin}
\externaldocument[quot-]{quot}
\externaldocument[stacks-properties-]{stacks-properties}
\externaldocument[stacks-morphisms-]{stacks-morphisms}
\externaldocument[stacks-limits-]{stacks-limits}
\externaldocument[stacks-cohomology-]{stacks-cohomology}
\externaldocument[stacks-perfect-]{stacks-perfect}
\externaldocument[stacks-introduction-]{stacks-introduction}
\externaldocument[stacks-more-morphisms-]{stacks-more-morphisms}
\externaldocument[stacks-geometry-]{stacks-geometry}
\externaldocument[moduli-]{moduli}
\externaldocument[moduli-curves-]{moduli-curves}
\externaldocument[examples-]{examples}
\externaldocument[exercises-]{exercises}
\externaldocument[guide-]{guide}
\externaldocument[desirables-]{desirables}
\externaldocument[coding-]{coding}
\externaldocument[obsolete-]{obsolete}
\externaldocument[fdl-]{fdl}
\externaldocument[index-]{index}

% Theorem environments.
%
\theoremstyle{plain}
\newtheorem{theorem}[subsection]{Theorem}
\newtheorem{proposition}[subsection]{Proposition}
\newtheorem{lemma}[subsection]{Lemma}

\theoremstyle{definition}
\newtheorem{definition}[subsection]{Definition}
\newtheorem{example}[subsection]{Example}
\newtheorem{exercise}[subsection]{Exercise}
\newtheorem{situation}[subsection]{Situation}

\theoremstyle{remark}
\newtheorem{remark}[subsection]{Remark}
\newtheorem{remarks}[subsection]{Remarks}

\numberwithin{equation}{subsection}

% Macros
%
\def\lim{\mathop{\mathrm{lim}}\nolimits}
\def\colim{\mathop{\mathrm{colim}}\nolimits}
\def\Spec{\mathop{\mathrm{Spec}}}
\def\Hom{\mathop{\mathrm{Hom}}\nolimits}
\def\Ext{\mathop{\mathrm{Ext}}\nolimits}
\def\SheafHom{\mathop{\mathcal{H}\!\mathit{om}}\nolimits}
\def\SheafExt{\mathop{\mathcal{E}\!\mathit{xt}}\nolimits}
\def\Sch{\mathit{Sch}}
\def\Mor{\mathop{\mathrm{Mor}}\nolimits}
\def\Ob{\mathop{\mathrm{Ob}}\nolimits}
\def\Sh{\mathop{\mathit{Sh}}\nolimits}
\def\NL{\mathop{N\!L}\nolimits}
\def\CH{\mathop{\mathrm{CH}}\nolimits}
\def\proetale{{pro\text{-}\acute{e}tale}}
\def\etale{{\acute{e}tale}}
\def\QCoh{\mathit{QCoh}}
\def\Ker{\mathop{\mathrm{Ker}}}
\def\Im{\mathop{\mathrm{Im}}}
\def\Coker{\mathop{\mathrm{Coker}}}
\def\Coim{\mathop{\mathrm{Coim}}}

% Boxtimes
%
\DeclareMathSymbol{\boxtimes}{\mathbin}{AMSa}{"02}

%
% Macros for moduli stacks/spaces
%
\def\QCohstack{\mathcal{QC}\!\mathit{oh}}
\def\Cohstack{\mathcal{C}\!\mathit{oh}}
\def\Spacesstack{\mathcal{S}\!\mathit{paces}}
\def\Quotfunctor{\mathrm{Quot}}
\def\Hilbfunctor{\mathrm{Hilb}}
\def\Curvesstack{\mathcal{C}\!\mathit{urves}}
\def\Polarizedstack{\mathcal{P}\!\mathit{olarized}}
\def\Complexesstack{\mathcal{C}\!\mathit{omplexes}}
% \Pic is the operator that assigns to X its picard group, usage \Pic(X)
% \Picardstack_{X/B} denotes the Picard stack of X over B
% \Picardfunctor_{X/B} denotes the Picard functor of X over B
\def\Pic{\mathop{\mathrm{Pic}}\nolimits}
\def\Picardstack{\mathcal{P}\!\mathit{ic}}
\def\Picardfunctor{\mathrm{Pic}}
\def\Deformationcategory{\mathcal{D}\!\mathit{ef}}
\usepackage{cleveref}
\crefname{originalequation}{original equation}{original equations}
\Crefname{originalequation}{Original equation}{Original equations}
\makeatletter
\newlabel{equation-finite-conditions}{{04PC}{source}{}{sourcepointer-1}{}}
\newlabel{equation-finite-conditions@cref}{{[originalequation][0][]04PC}{[1][0][]source}{}{}{}}
\makeatother
\crefname{originallemma}{original lemma}{original lemmas}
\Crefname{originallemma}{Original lemma}{Original lemmas}
\makeatletter
\newlabel{lemma-finite-criterion-etale}{{04QE}{source}{}{sourcepointer-2}{}}
\newlabel{lemma-finite-criterion-etale@cref}{{[originallemma][0][]04QE}{[1][0][]source}{}{}{}}
\makeatother
\crefname{originallemma}{original lemma}{original lemmas}
\Crefname{originallemma}{Original lemma}{Original lemmas}
\makeatletter
\newlabel{lemma-finite-diagonal}{{04PK}{source}{}{sourcepointer-3}{}}
\newlabel{lemma-finite-diagonal@cref}{{[originallemma][0][]04PK}{[1][0][]source}{}{}{}}
\makeatother
\crefname{originallemma}{original lemma}{original lemmas}
\Crefname{originallemma}{Original lemma}{Original lemmas}
\makeatletter
\newlabel{lemma-finite-pullback}{{04QF}{source}{}{sourcepointer-4}{}}
\newlabel{lemma-finite-pullback@cref}{{[originallemma][0][]04QF}{[1][0][]source}{}{}{}}
\makeatother
\crefname{originalreference}{original reference}{original references}
\Crefname{originalreference}{Original reference}{Original references}
\makeatletter
\newlabel{more-morphisms-lemma-etale-splits-off-quasi-finite-part-technical-variant}{{02LO}{source}{}{sourcepointer-5}{}}
\newlabel{more-morphisms-lemma-etale-splits-off-quasi-finite-part-technical-variant@cref}{{[originalreference][0][]02LO}{[1][0][]source}{}{}{}}
\makeatother
\crefname{originalreference}{original reference}{original references}
\Crefname{originalreference}{Original reference}{Original references}
\makeatletter
\newlabel{morphisms-lemma-composition-etale}{{02GN}{source}{}{sourcepointer-6}{}}
\newlabel{morphisms-lemma-composition-etale@cref}{{[originalreference][0][]02GN}{[1][0][]source}{}{}{}}
\makeatother
\crefname{originalreference}{original reference}{original references}
\Crefname{originalreference}{Original reference}{Original references}
\makeatletter
\newlabel{morphisms-lemma-composition-surjective}{{01S0}{source}{}{sourcepointer-7}{}}
\newlabel{morphisms-lemma-composition-surjective@cref}{{[originalreference][0][]01S0}{[1][0][]source}{}{}{}}
\makeatother
\crefname{originalreference}{original reference}{original references}
\Crefname{originalreference}{Original reference}{Original references}
\makeatletter
\newlabel{spaces-lemma-base-change-representable-transformations-property}{{02WL}{source}{}{sourcepointer-8}{}}
\newlabel{spaces-lemma-base-change-representable-transformations-property@cref}{{[originalreference][0][]02WL}{[1][0][]source}{}{}{}}
\makeatother
\crefname{originalreference}{original reference}{original references}
\Crefname{originalreference}{Original reference}{Original references}
\makeatletter
\newlabel{spaces-lemma-composition-representable-transformations}{{02WA}{source}{}{sourcepointer-9}{}}
\newlabel{spaces-lemma-composition-representable-transformations@cref}{{[originalreference][0][]02WA}{[1][0][]source}{}{}{}}
\makeatother
\crefname{originalreference}{original reference}{original references}
\Crefname{originalreference}{Original reference}{Original references}
\makeatletter
\newlabel{spaces-lemma-composition-representable-transformations-property}{{02WK}{source}{}{sourcepointer-10}{}}
\newlabel{spaces-lemma-composition-representable-transformations-property@cref}{{[originalreference][0][]02WK}{[1][0][]source}{}{}{}}
\makeatother
\crefname{originalreference}{original reference}{original references}
\Crefname{originalreference}{Original reference}{Original references}
\makeatletter
\newlabel{spaces-lemma-representable-diagonal}{{025W}{source}{}{sourcepointer-11}{}}
\newlabel{spaces-lemma-representable-diagonal@cref}{{[originalreference][0][]025W}{[1][0][]source}{}{}{}}
\makeatother
\crefname{originalreference}{original reference}{original references}
\Crefname{originalreference}{Original reference}{Original references}
\makeatletter
\newlabel{spaces-morphisms-lemma-base-change-quasi-finite}{{03XL}{source}{}{sourcepointer-12}{}}
\newlabel{spaces-morphisms-lemma-base-change-quasi-finite@cref}{{[originalreference][0][]03XL}{[1][0][]source}{}{}{}}
\makeatother
\crefname{originalreference}{original reference}{original references}
\Crefname{originalreference}{Original reference}{Original references}
\makeatletter
\newlabel{spaces-morphisms-lemma-base-change-separated}{{03KL}{source}{}{sourcepointer-13}{}}
\newlabel{spaces-morphisms-lemma-base-change-separated@cref}{{[originalreference][0][]03KL}{[1][0][]source}{}{}{}}
\makeatother
\crefname{originalreference}{original reference}{original references}
\Crefname{originalreference}{Original reference}{Original references}
\makeatletter
\newlabel{spaces-morphisms-proposition-locally-quasi-finite-separated-over-scheme}{{03XX}{source}{}{sourcepointer-14}{}}
\newlabel{spaces-morphisms-proposition-locally-quasi-finite-separated-over-scheme@cref}{{[originalreference][0][]03XX}{[1][0][]source}{}{}{}}
\makeatother
\crefname{originalreference}{original reference}{original references}
\Crefname{originalreference}{Original reference}{Original references}
\makeatletter
\newlabel{spaces-remark-warning}{{02YN}{source}{}{sourcepointer-15}{}}
\newlabel{spaces-remark-warning@cref}{{[originalreference][0][]02YN}{[1][0][]source}{}{}{}}
\makeatother
\title{The finite part of a separated locally quasi-finite map is covered by a scheme etale over Y}
\author{The Stacks Project Authors}
\date{stacks-project:spaces-more-groupoids.tex:local-snapshot-20261012; GNU FDL 1.2 or later}
\begin{document}
\maketitle
\noindent\textit{Editorial note.} No omission markers in-span; two commutative diagrams are displayed and used. Dependence on the uncarded base-change lemma is declared as external.\par\medskip
\noindent\textit{Allowed standard prerequisites (editorial).}\begin{itemize}
\item Etale neighbourhood splitting off the quasi-finite part. More on Morphisms, lemma etale-splits-off-quasi-finite-part-technical-variant (cross-chapter, named). Supplies the decomposition U x\_S X = W amal coprod V\_\{i,j\} with purely inseparable residue fields at the v\_\{i,j\}.
\item Base change of the finite part. lemma-finite-pullback (same chapter, NOT carded, declared external). Identifies (V x\_Y X/V)\_fin as a base change, surjective and etale.
\item Etale criterion and the diagonal. lemma-finite-criterion-etale (carded L26654) and lemma-finite-diagonal (carded L26653). Etaleness of the constructed morphisms and the representability used throughout.
\item Composition of representable transformation properties. Spaces and Morphisms lemma block on representable/surjective/etale transformations (cross-chapter, named). Justifies replacing Y by the cover V.
\end{itemize}\medskip
\ifdefined\hypertarget\hypertarget{sourcepointer-1}{}\fi
\noindent\textit{Source reference.} Original equation 04PC: sources/stacks-spaces-more-groupoids/pinned-original-source.tex, source lines 1519--1519. Navigation pointer only: the counted proof cites this statement by \textbackslash{}ref as a standing input; the deduction the counted proof has to make is carried out inside the counted span, and no part of that deduction is delegated to this pointer.\par\medskip
\ifdefined\hypertarget\hypertarget{sourcepointer-2}{}\fi
\noindent\textit{Source reference.} Original lemma 04QE: sources/stacks-spaces-more-groupoids/pinned-original-source.tex, source lines 1790--1790. Navigation pointer only: the counted proof cites this statement by \textbackslash{}ref as a standing input; the deduction the counted proof has to make is carried out inside the counted span, and no part of that deduction is delegated to this pointer.\par\medskip
\ifdefined\hypertarget\hypertarget{sourcepointer-3}{}\fi
\noindent\textit{Source reference.} Original lemma 04PK: sources/stacks-spaces-more-groupoids/pinned-original-source.tex, source lines 1744--1744. Navigation pointer only: the counted proof cites this statement by \textbackslash{}ref as a standing input; the deduction the counted proof has to make is carried out inside the counted span, and no part of that deduction is delegated to this pointer.\par\medskip
\ifdefined\hypertarget\hypertarget{sourcepointer-4}{}\fi
\noindent\textit{Source reference.} Original lemma 04QF: sources/stacks-spaces-more-groupoids/pinned-original-source.tex, source lines 1822--1822. Navigation pointer only: the counted proof cites this statement by \textbackslash{}ref as a standing input; the deduction the counted proof has to make is carried out inside the counted span, and no part of that deduction is delegated to this pointer.\par\medskip
\ifdefined\hypertarget\hypertarget{sourcepointer-5}{}\fi
\noindent\textit{Source reference.} Original reference 02LO: focused-seeds/geometry/S9-smg-finite-surjective-etale-cover/ref-tags.txt, source lines 5--5. Navigation pointer only: the counted proof cites this statement by \textbackslash{}ref as a standing input; the deduction the counted proof has to make is carried out inside the counted span, and no part of that deduction is delegated to this pointer.\par\medskip
\ifdefined\hypertarget\hypertarget{sourcepointer-6}{}\fi
\noindent\textit{Source reference.} Original reference 02GN: focused-seeds/geometry/S9-smg-finite-surjective-etale-cover/ref-tags.txt, source lines 6--6. Navigation pointer only: the counted proof cites this statement by \textbackslash{}ref as a standing input; the deduction the counted proof has to make is carried out inside the counted span, and no part of that deduction is delegated to this pointer.\par\medskip
\ifdefined\hypertarget\hypertarget{sourcepointer-7}{}\fi
\noindent\textit{Source reference.} Original reference 01S0: focused-seeds/geometry/S9-smg-finite-surjective-etale-cover/ref-tags.txt, source lines 7--7. Navigation pointer only: the counted proof cites this statement by \textbackslash{}ref as a standing input; the deduction the counted proof has to make is carried out inside the counted span, and no part of that deduction is delegated to this pointer.\par\medskip
\ifdefined\hypertarget\hypertarget{sourcepointer-8}{}\fi
\noindent\textit{Source reference.} Original reference 02WL: focused-seeds/geometry/S9-smg-finite-surjective-etale-cover/ref-tags.txt, source lines 8--8. Navigation pointer only: the counted proof cites this statement by \textbackslash{}ref as a standing input; the deduction the counted proof has to make is carried out inside the counted span, and no part of that deduction is delegated to this pointer.\par\medskip
\ifdefined\hypertarget\hypertarget{sourcepointer-9}{}\fi
\noindent\textit{Source reference.} Original reference 02WA: focused-seeds/geometry/S9-smg-finite-surjective-etale-cover/ref-tags.txt, source lines 9--9. Navigation pointer only: the counted proof cites this statement by \textbackslash{}ref as a standing input; the deduction the counted proof has to make is carried out inside the counted span, and no part of that deduction is delegated to this pointer.\par\medskip
\ifdefined\hypertarget\hypertarget{sourcepointer-10}{}\fi
\noindent\textit{Source reference.} Original reference 02WK: focused-seeds/geometry/S9-smg-finite-surjective-etale-cover/ref-tags.txt, source lines 10--10. Navigation pointer only: the counted proof cites this statement by \textbackslash{}ref as a standing input; the deduction the counted proof has to make is carried out inside the counted span, and no part of that deduction is delegated to this pointer.\par\medskip
\ifdefined\hypertarget\hypertarget{sourcepointer-11}{}\fi
\noindent\textit{Source reference.} Original reference 025W: focused-seeds/geometry/S9-smg-finite-surjective-etale-cover/ref-tags.txt, source lines 11--11. Navigation pointer only: the counted proof cites this statement by \textbackslash{}ref as a standing input; the deduction the counted proof has to make is carried out inside the counted span, and no part of that deduction is delegated to this pointer.\par\medskip
\ifdefined\hypertarget\hypertarget{sourcepointer-12}{}\fi
\noindent\textit{Source reference.} Original reference 03XL: focused-seeds/geometry/S9-smg-finite-surjective-etale-cover/ref-tags.txt, source lines 12--12. Navigation pointer only: the counted proof cites this statement by \textbackslash{}ref as a standing input; the deduction the counted proof has to make is carried out inside the counted span, and no part of that deduction is delegated to this pointer.\par\medskip
\ifdefined\hypertarget\hypertarget{sourcepointer-13}{}\fi
\noindent\textit{Source reference.} Original reference 03KL: focused-seeds/geometry/S9-smg-finite-surjective-etale-cover/ref-tags.txt, source lines 13--13. Navigation pointer only: the counted proof cites this statement by \textbackslash{}ref as a standing input; the deduction the counted proof has to make is carried out inside the counted span, and no part of that deduction is delegated to this pointer.\par\medskip
\ifdefined\hypertarget\hypertarget{sourcepointer-14}{}\fi
\noindent\textit{Source reference.} Original reference 03XX: focused-seeds/geometry/S9-smg-finite-surjective-etale-cover/ref-tags.txt, source lines 14--14. Navigation pointer only: the counted proof cites this statement by \textbackslash{}ref as a standing input; the deduction the counted proof has to make is carried out inside the counted span, and no part of that deduction is delegated to this pointer.\par\medskip
\ifdefined\hypertarget\hypertarget{sourcepointer-15}{}\fi
\noindent\textit{Source reference.} Original reference 02YN: focused-seeds/geometry/S9-smg-finite-surjective-etale-cover/ref-tags.txt, source lines 15--15. Navigation pointer only: the counted proof cites this statement by \textbackslash{}ref as a standing input; the deduction the counted proof has to make is carried out inside the counted span, and no part of that deduction is delegated to this pointer.\par\medskip
\begin{lemma}
\label{lemma-finite-surjective-etale-cover}
Let $S$ be a scheme.
Let $f : X \to Y$ be a morphism of algebraic spaces over $S$.
If $f$ is separated and locally quasi-finite, then there exists a
scheme $U$ \'etale over $Y$ and a surjective \'etale morphism
$U \to (X/Y)_{fin}$ over $Y$.
\end{lemma}

\begin{proof}
Note that the assertion makes sense by the result of
Lemma \ref{lemma-finite-diagonal}
on the diagonal of $(X/Y)_{fin}$, see
Spaces, Lemma \ref{spaces-lemma-representable-diagonal}.
Let $V$ be a scheme and let $V \to Y$ be a surjective \'etale morphism. By
Lemma \ref{lemma-finite-pullback}
the morphism $(V \times_Y X/V)_{fin} \to (X/Y)_{fin}$ is
a base change of the map $V \to Y$ and hence is surjective and \'etale, see
Spaces,
Lemma \ref{spaces-lemma-base-change-representable-transformations-property}.
Hence it suffices to prove the lemma for $(V \times_Y X/V)_{fin}$.
(Here we implicitly use that the composition of representable, surjective, and
\'etale transformations of functors is again representable, surjective, and
\'etale, see
Spaces, Lemmas \ref{spaces-lemma-composition-representable-transformations} and
\ref{spaces-lemma-composition-representable-transformations-property}, and
Morphisms, Lemmas \ref{morphisms-lemma-composition-surjective} and
\ref{morphisms-lemma-composition-etale}.)
Note that the properties of being separated and locally quasi-finite
are preserved under base change, see
Morphisms of Spaces,
Lemmas \ref{spaces-morphisms-lemma-base-change-separated} and
\ref{spaces-morphisms-lemma-base-change-quasi-finite}.
Hence $V \times_Y X \to V$ is separated and locally quasi-finite as well,
and by
Morphisms of Spaces, Proposition
\ref{spaces-morphisms-proposition-locally-quasi-finite-separated-over-scheme}
we see that $V \times_Y X$ is a scheme as well.
Thus we may assume that $f : X \to Y$ is a separated and locally quasi-finite
morphism of schemes.

\medskip\noindent
Pick a point $y \in Y$. Pick $x_1, \ldots, x_n \in X$ points
lying over $y$. Pick an \'etale neighbourhood $a : (U, u) \to (Y, y)$ and a
decomposition
$$
U \times_S X =
W \amalg
\ \coprod\nolimits_{i = 1, \ldots, n}
\ \coprod\nolimits_{j = 1, \ldots, m_j}
V_{i, j}
$$
as in
More on Morphisms, Lemma
\ref{more-morphisms-lemma-etale-splits-off-quasi-finite-part-technical-variant}.
Pick any subset
$$
I \subset \{(i, j) \mid 1 \leq i \leq n, \ 1 \leq j \leq m_i\}.
$$
Given these choices we obtain a pair $(a, Z)$ with
$Z = \bigcup_{(i, j) \in I} V_{i, j}$
which satisfies conditions \ref{equation-finite-conditions}. In other words
we obtain a morphism $U \to (X/Y)_{fin}$. The construction of this morphism
depends on all the things we picked above, so we should really write
$$
U(y, n, x_1, \ldots, x_n, a, I) \longrightarrow (X/Y)_{fin}
$$
This morphism is \'etale by Lemma \ref{lemma-finite-criterion-etale}.

\medskip\noindent
Claim: The disjoint union of all of these is surjective onto $(X/Y)_{fin}$.
It is clear that if the claim holds, then the lemma is true.

\medskip\noindent
To show surjectivity we have to show the following (see
Spaces, Remark \ref{spaces-remark-warning}): Given a scheme $T$ over
$S$, a point $t \in T$, and a map $T \to (X/Y)_{fin}$ we can find a datum
$(y, n, x_1, \ldots, x_n, a, I)$ as above such that
$t$ is in the image of the projection map
$$
U(y, n, x_1, \ldots, x_n, a, I) \times_{(X/Y)_{fin}} T \longrightarrow T.
$$
To prove this we may clearly replace $T$ by
$\Spec(\overline{\kappa(t)})$ and
$T \to (X/Y)_{fin}$ by the composition
$\Spec(\overline{\kappa(t)}) \to T \to (X/Y)_{fin}$.
In other words, we may assume that $T$ is
the spectrum of an algebraically closed field.

\medskip\noindent
Let $T = \Spec(k)$ be the spectrum of an algebraically closed
field $k$. The morphism $T \to (X/Y)_{fin}$ is given by a pair
$(T \to Y, Z)$ satisfying conditions \ref{equation-finite-conditions}.
Here is a picture:
$$
\xymatrix{
& Z \ar[d] \ar[r] & X \ar[d] \\
\Spec(k) \ar@{=}[r] & T \ar[r] & Y
}
$$
Let $y \in Y$ be the image point of $T \to Y$.
Since $Z$ is finite over $k$ it has finitely many points.
Thus there exist finitely many points $x_1, \ldots, x_n \in X$
such that the image of $Z$ in $X$ is contained in $\{x_1, \ldots, x_n\}$.
Choose $a : (U, u) \to (Y, y)$ adapted to $y$ and $x_1, \ldots, x_n$ as
above, which gives the diagram
$$
\xymatrix{
W \amalg
\ \coprod\nolimits_{i = 1, \ldots, n}
\ \coprod\nolimits_{j = 1, \ldots, m_j}
V_{i, j} \ar[d] \ar[r] &
X \ar[d] \\
U \ar[r] & Y.
}
$$
Since $k$ is algebraically closed and
$\kappa(y) \subset \kappa(u)$ is finite separable
we may factor the morphism
$T = \Spec(k) \to Y$ through the morphism
$u = \Spec(\kappa(u)) \to \Spec(\kappa(y)) = y \subset Y$.
With this choice we obtain the commutative diagram:
$$
\xymatrix{
Z \ar[d] \ar[r] &
W \amalg
\ \coprod\nolimits_{i = 1, \ldots, n}
\ \coprod\nolimits_{j = 1, \ldots, m_j}
V_{i, j} \ar[d] \ar[r] &
X \ar[d] \\
\Spec(k) \ar[r] &
U \ar[r] & Y
}
$$
We know that the image of the left upper arrow ends up in
$\coprod V_{i, j}$. Recall also that $Z$ is an open subscheme
of $\Spec(k) \times_Y X$ by definition of $(X/Y)_{fin}$
and that the right hand square is a fibre product square.
Thus we see that
$$
Z \subset
\coprod\nolimits_{i = 1, \ldots, n}\ \coprod\nolimits_{j = 1, \ldots, m_j}
\Spec(k) \times_U V_{i, j}
$$
is an open subscheme. By construction (see
More on Morphisms, Lemma
\ref{more-morphisms-lemma-etale-splits-off-quasi-finite-part-technical-variant})
each $V_{i, j}$ has a unique point $v_{i, j}$ lying over $u$
with purely inseparable residue field extension
$\kappa(v_{i, j})/\kappa(u)$. Hence each
scheme $\Spec(k) \times_U V_{i, j}$ has exactly one
point. Thus we see that
$$
Z = \coprod\nolimits_{(i, j) \in I} \Spec(k) \times_U V_{i, j}
$$
for a unique subset
$I \subset \{(i, j) \mid 1 \leq i \leq n, \ 1 \leq j \leq m_i\}$.
Unwinding the definitions this shows that
$$
U(y, n, x_1, \ldots, x_n, a, I) \times_{(X/Y)_{fin}} T
$$
with $I$ as found above is nonempty as desired.
\end{proof}

\end{document}
